\chapter{О чём это руководство}
\label{ch:intro}

RPG Master --- платформа для проведения настольных ролевых игр через интернет. Мастер готовит сценарий заранее, а во время игры управляет сценами, картой, музыкой и бросками; игроки видят только то, что мастер им открыл.

\paragraph{Для кого}
Руководство рассчитано на трёх читателей, и каждому достаточно своей части:

\begin{itemize}
\item \textbf{Игроку} --- главы \ref{ch:getting-started} и \ref{ch:player}. Этого хватит, чтобы войти, получить персонажа и сыграть.
\item \textbf{Мастеру} --- дополнительно главы \ref{ch:concepts} и \ref{ch:master}. Глава \ref{ch:concepts} объясняет, чем сценарий отличается от запущенного мира; без неё легко запутаться.
\item \textbf{Зрителю} --- глава \ref{ch:observer}. Регистрация не нужна, достаточно ссылки от мастера.
\end{itemize}

\paragraph{Как устроено}
Первые главы читаются подряд и ведут от регистрации до конца первой игры. Дальше идут справочники: сущности сценария (\ref{ch:objects}), экраны сайта (\ref{ch:navigation}) и типовые действия (\ref{ch:standard-actions}) --- в них заглядывают по мере надобности, а не читают целиком. В конце --- глава \ref{ch:troubleshooting} с решениями частых проблем.

\paragraph{Условные обозначения}

\begin{itemize}
\item \texttt{/scenarios} --- адрес страницы. Подставляется к адресу сайта: если сайт открывается как \texttt{example.ru}, то страница будет \texttt{example.ru/scenarios}.
\item ``Сценарии'' --- надпись на кнопке или вкладке, ровно как она выглядит на экране.
\item \emph{Сценарий}, \emph{подход} --- термины платформы. Все они собраны в главе \ref{ch:concepts}.
\end{itemize}

\paragraph{Чего в руководстве нет}
Здесь не объясняются сами правила ролевых игр. Платформа поддерживает несколько систем (Dungeon World, Gumshoe, Blades in the Dark и другие), но учит играть по ним не она --- нужны книги правил соответствующей системы. Руководство описывает только то, как пользоваться сайтом.

Технические подробности --- устройство плагинов, формат данных, протокол обмена --- вынесены в отдельный документ «Техническая документация» и здесь не рассматриваются.
